\section{Conclusion and Future Work}
\label{sec:conclusion}

This paper presented ContrastiveXAI-FS, an unsupervised representation-aware feature-selection and cluster-interpretation pipeline for IIoT intrusion detection. The method combines group-aware VICReg representation learning (C1), bidirectional Silhouette-based feature refinement in the learned latent space (C2), and SHAP-based cluster profiling through a Random Forest proxy (C3). Class labels are not used during representation learning, feature selection, or cluster construction.

\rev{On the DataSense CIC IIoT 2025 dataset, with every stage fitted inside the training partition of each fold, ContrastiveXAI-FS retains on average 59 of the 71 numeric features and achieves Random Forest macro-F1 scores of 0.936 for binary detection and 0.908 for eight-class classification. Its performance is statistically equivalent to the complete feature set within a margin of $\pm$0.01 macro-F1, consistently across three independent master seeds, and it significantly outperforms MCFS at its natural budget and the supervised DataSense-17 subset. At identical budgets, evaluating candidate subsets through the VICReg latent space yields higher performance and more stable subsets than the same search in the raw feature space. The same pattern holds on N-BaIoT, where the method is equivalent to the complete feature set and superior to MCFS, and the method was applied without tuning to WUSTL-IIoT-2021, where it is competitive in binary detection.}

\rev{The explainability stage is supported by direct evidence: the latent clusters are reproducible across folds and master seeds (adjusted Rand index 0.94), the Random Forest proxy reproduces them with 0.97 balanced accuracy on unseen test folds, and a quantitative clarity index identifies the clusters whose profiles are both faithful and concentrated. The resulting profiles, dominated by TCP termination flags, TTL levels, TTL variability, and payload sizes, connect the unlabeled cluster structure back to measurements that security analysts can inspect.}

\rev{The method is best suited to representation-aware refinement at moderate to large budgets, whereas compact classical selectors remain preferable under very small feature budgets. The main results use sample-level folds and therefore measure within-dataset generalization; device- and execution-grouped evaluations show that binary detection is preserved for held-out attack executions and remains above 0.90 for unseen devices, while multi-class recognition of categories recorded on few devices is harder for every configuration.}

Future work could address four directions. First, budget-aware variants of the latent search could extend the operating region to very small feature budgets. Second, hybrid strategies could combine VICReg-based preselection with raw-space refinement for edge deployment. Third, online and incremental versions could be studied for environments in which devices are added over time and traffic distributions change. \rev{Fourth, evaluation on additional industrial datasets with device- and time-grouped splits and analyst-in-the-loop validation of the behavioral profiles would further characterize deployment-level generalization.}
